\documentclass[11pt,a4paper]{article}

\usepackage[margin=18mm]{geometry}
\usepackage[T1]{fontenc}
\usepackage[utf8]{inputenc}
\usepackage{amsmath}
\usepackage{array}
\usepackage{booktabs}
\usepackage{caption}
\usepackage{enumitem}
\usepackage{fancyhdr}
\usepackage{float}
\usepackage{graphicx}
\usepackage{hyperref}
\usepackage{microtype}
\usepackage{tabularx}
\usepackage[table]{xcolor}

% -----------------------------------------------------------------------------
% Replace these values before submitting.
% The assignment requires the website URL immediately below the assignment
% heading and requires the link to open without requesting access permission.
% -----------------------------------------------------------------------------
\newcommand{\websiteurl}{https://public.emirhanaltunel.com/homeworks/cse513-1/}
\newcommand{\studentname}{Emirhan Altunel}
\newcommand{\studentid}{26420101003}

\definecolor{ink}{HTML}{12203A}
\definecolor{accent}{HTML}{2457F5}
\definecolor{softblue}{HTML}{EEF2FF}
\definecolor{softgray}{HTML}{F3F6FB}
\definecolor{success}{HTML}{207C4B}

\hypersetup{
  colorlinks=true,
  urlcolor=accent,
  linkcolor=accent,
  pdftitle={CSE513 Assignment 1 -- CPU Scheduling Simulator},
  pdfauthor={\studentname}
}

\pagestyle{fancy}
\fancyhf{}
\lhead{CSE513 Assignment 1}
\rhead{CPU Scheduling Simulator}
\cfoot{\thepage}
\setlength{\headheight}{14pt}
\setlength{\parindent}{0pt}
\setlength{\parskip}{5pt}
\renewcommand{\arraystretch}{1.12}
\setlist{nosep,leftmargin=1.5em}

\newcommand{\timeline}[1]{%
  \begin{center}
  \colorbox{softgray}{\parbox{0.94\linewidth}{\ttfamily\small #1}}
  \end{center}
}

\newcommand{\exactmatch}{\textcolor{success}{\textbf{Exact match; no discrepancy.}}}

% Add a screenshot by placing it in report/screenshots/ with the requested name.
% The second path allows compilation from the project root as well as report/.
\newcommand{\simshot}[2]{%
  \begin{figure}[H]
    \centering
    \IfFileExists{screenshots/#1}{%
      \includegraphics[width=\linewidth,height=0.34\textheight,keepaspectratio]{screenshots/#1}%
    }{%
      \IfFileExists{report/screenshots/#1}{%
        \includegraphics[width=\linewidth,height=0.34\textheight,keepaspectratio]{report/screenshots/#1}%
      }{%
        \fbox{\begin{minipage}[c][3.4cm][c]{0.91\linewidth}
          \centering
          \textbf{Screenshot placeholder}\\[4pt]
          Add \texttt{#1} to \texttt{report/screenshots/}, then recompile.
        \end{minipage}}%
      }%
    }
    \caption{#2}
  \end{figure}
}

\begin{document}

% First-page ordering follows the assignment instructions exactly.
{\large\bfseries\color{accent} CSE513 Assignment 1}\par
\vspace{3pt}
\textbf{Working website: }\href{\websiteurl}{\websiteurl}\par
\vspace{16mm}

{\Huge\bfseries\color{ink} CPU Scheduling Simulator}\par
\vspace{6mm}
{\large A deterministic browser-based simulator for FCFS, SJF, SRTF, and Round Robin.}\par
\vspace{15mm}

\begin{tabular}{@{}ll@{}}
  \textbf{Student:}    & \studentname \\
  \textbf{Student ID:} & \studentid   \\
\end{tabular}

\vfill
\colorbox{softblue}{\parbox{0.92\linewidth}{
    \textbf{Report scope.} This report presents the scheduling model, eight required
    algorithm--workload examples, hand calculations, simulator comparisons, three
    additional test cases, invariant checks, comparative analysis, implementation
    notes, and an AI-assistance statement.
  }}
\vfill
\begin{center}
  \begin{tabular}{cccc}
    \textbf{4 algorithms} & \textbf{8 worked examples} & \textbf{3 custom tests} & \textbf{9 automated tests}
  \end{tabular}
\end{center}
\newpage

\tableofcontents
\newpage

\section{Model and calculation method}

The simulator models one work-conserving CPU, one CPU burst per process, no I/O,
and zero-time context switches. Arrival times are nonnegative integers; CPU bursts
and the Round Robin quantum are positive integers; and process IDs are unique.
Intervals are half-open, so $[2,5)$ supplies three CPU units.

For every process, the metrics are calculated as follows:

\begin{center}
  \colorbox{softblue}{\parbox{0.75\linewidth}{
      \begin{align*}
        T & = C-A  &  & \text{(turnaround time)}, \\
        W & = T-B  &  & \text{(waiting time)},    \\
        R & = FS-A &  & \text{(response time)}.
      \end{align*}
    }}
\end{center}

Ties prefer the earlier arrival time and then the earlier position in the input
list. SRTF preempts only when another ready process has a \emph{strictly} shorter
remaining time. At a Round Robin quantum boundary, arrivals at that exact time
enter the ready queue before the unfinished running process is re-enqueued.

\subsection{Required input workloads}

\begin{table}[H]
  \centering
  \caption{Workloads supplied in the assignment.}
  \begin{tabular}{lrr@{\qquad}lrr}
    \toprule
    \multicolumn{3}{c}{\textbf{Workload A}} & \multicolumn{3}{c}{\textbf{Workload B}}                                     \\
    Process                                 & Arrival                                 & Burst & Process & Arrival & Burst \\
    \midrule
    P1                                      & 0                                       & 8     & P1      & 2       & 3     \\
    P2                                      & 1                                       & 4     & P2      & 2       & 1     \\
    P3                                      & 2                                       & 2     & P3      & 8       & 2     \\
    P4                                      & 4                                       & 1     & P4      & 9       & 1     \\
    \bottomrule
  \end{tabular}
\end{table}

\section{Worked examples: Workload A}

\subsection{FCFS on Workload A}

\textbf{Input:} P1=$(0,8)$, P2=$(1,4)$, P3=$(2,2)$, P4=$(4,1)$, where each pair
is (arrival, burst).

\textbf{Hand decision trace.} P1 is the only process ready at time 0 and runs to
completion. At time 8, P2, P3, and P4 are ready. FCFS chooses them by arrival
order: P2, P3, then P4.

\timeline{0--8 P1 \;|\; 8--12 P2 \;|\; 12--14 P3 \;|\; 14--15 P4}

\begin{table}[H]
  \centering\small
  \caption{FCFS Workload A: hand calculation and simulator output (identical).}
  \begin{tabular}{lrrrrrrr}
    \toprule
    Process                              & A              & B             & FS            & C  & T  & W  & R  \\
    \midrule
    P1                                   & 0              & 8             & 0             & 8  & 8  & 0  & 0  \\
    P2                                   & 1              & 4             & 8             & 12 & 11 & 7  & 7  \\
    P3                                   & 2              & 2             & 12            & 14 & 12 & 10 & 10 \\
    P4                                   & 4              & 1             & 14            & 15 & 11 & 10 & 10 \\
    \midrule
    \multicolumn{5}{r}{\textbf{Average}} & \textbf{10.50} & \textbf{6.75} & \textbf{6.75}                     \\
    \bottomrule
  \end{tabular}
\end{table}

\textbf{Simulator comparison.} The simulator produced the same timeline and every
per-process value in the table. \exactmatch

\simshot{a-fcfs.png}{Simulator output for FCFS on Workload A.}

\subsection{SJF on Workload A}

\textbf{Input:} P1=$(0,8)$, P2=$(1,4)$, P3=$(2,2)$, P4=$(4,1)$.

\textbf{Hand decision trace.} Non-preemptive P1 must finish. At time 8 all other
jobs are ready, so ascending burst length gives P4 (1), P3 (2), and P2 (4).

\timeline{0--8 P1 \;|\; 8--9 P4 \;|\; 9--11 P3 \;|\; 11--15 P2}

\begin{table}[H]
  \centering\small
  \caption{SJF Workload A: hand calculation and simulator output (identical).}
  \begin{tabular}{lrrrrrrr}
    \toprule
    Process                              & A             & B             & FS            & C  & T  & W  & R  \\
    \midrule
    P1                                   & 0             & 8             & 0             & 8  & 8  & 0  & 0  \\
    P2                                   & 1             & 4             & 11            & 15 & 14 & 10 & 10 \\
    P3                                   & 2             & 2             & 9             & 11 & 9  & 7  & 7  \\
    P4                                   & 4             & 1             & 8             & 9  & 5  & 4  & 4  \\
    \midrule
    \multicolumn{5}{r}{\textbf{Average}} & \textbf{9.00} & \textbf{5.25} & \textbf{5.25}                     \\
    \bottomrule
  \end{tabular}
\end{table}

\textbf{Simulator comparison.} The simulator produced the same timeline and every
per-process value in the table. \exactmatch

\simshot{a-sjf.png}{Simulator output for SJF on Workload A.}

\subsection{SRTF on Workload A}

\textbf{Input:} P1=$(0,8)$, P2=$(1,4)$, P3=$(2,2)$, P4=$(4,1)$.

\textbf{Hand decision trace.} P2 arrives at time 1 with burst 4, which is shorter
than P1's remaining 7, so it preempts P1. At time 2, P3's burst 2 is shorter than
P2's remaining 3. P3 finishes at time 4; P4 then has the shortest remaining time.
P2 resumes before P1.

\timeline{0--1 P1 \;|\; 1--2 P2 \;|\; 2--4 P3 \;|\; 4--5 P4 \;|\; 5--8 P2 \;|\; 8--15 P1}

\begin{table}[H]
  \centering\small
  \caption{SRTF Workload A: hand calculation and simulator output (identical).}
  \begin{tabular}{lrrrrrrr}
    \toprule
    Process                              & A             & B             & FS            & C  & T  & W & R \\
    \midrule
    P1                                   & 0             & 8             & 0             & 15 & 15 & 7 & 0 \\
    P2                                   & 1             & 4             & 1             & 8  & 7  & 3 & 0 \\
    P3                                   & 2             & 2             & 2             & 4  & 2  & 0 & 0 \\
    P4                                   & 4             & 1             & 4             & 5  & 1  & 0 & 0 \\
    \midrule
    \multicolumn{5}{r}{\textbf{Average}} & \textbf{6.25} & \textbf{2.50} & \textbf{0.00}                   \\
    \bottomrule
  \end{tabular}
\end{table}

\textbf{Simulator comparison.} The simulator produced the same timeline and every
per-process value in the table. \exactmatch

\simshot{a-srtf.png}{Simulator output for SRTF on Workload A.}

\subsection{Round Robin (q=2) on Workload A}

\textbf{Input:} P1=$(0,8)$, P2=$(1,4)$, P3=$(2,2)$, P4=$(4,1)$; quantum $q=2$.

\textbf{Hand decision trace.} Each process at the queue head receives at most two
units. P3's arrival at time 2 is enqueued before P1 is returned to the tail. P4
similarly enters at the time-4 boundary before P2 is re-enqueued. The final two
P1 slices remain separately visible even though they are consecutive.

\timeline{0--2 P1 \;|\; 2--4 P2 \;|\; 4--6 P3 \;|\; 6--8 P1 \;|\; 8--9 P4 \;|\; 9--11 P2 \;|\; 11--13 P1 \;|\; 13--15 P1}

\begin{table}[H]
  \centering\small
  \caption{Round Robin Workload A: hand calculation and simulator output (identical).}
  \begin{tabular}{lrrrrrrr}
    \toprule
    Process                              & A             & B             & FS            & C  & T  & W & R \\
    \midrule
    P1                                   & 0             & 8             & 0             & 15 & 15 & 7 & 0 \\
    P2                                   & 1             & 4             & 2             & 11 & 10 & 6 & 1 \\
    P3                                   & 2             & 2             & 4             & 6  & 4  & 2 & 2 \\
    P4                                   & 4             & 1             & 8             & 9  & 5  & 4 & 4 \\
    \midrule
    \multicolumn{5}{r}{\textbf{Average}} & \textbf{8.50} & \textbf{4.75} & \textbf{1.75}                   \\
    \bottomrule
  \end{tabular}
\end{table}

\textbf{Simulator comparison.} The simulator produced the same eight slices and
every per-process value in the table. \exactmatch

\simshot{a-rr.png}{Simulator output for Round Robin ($q=2$) on Workload A.}

\section{Worked examples: Workload B}

\subsection{FCFS on Workload B}

\textbf{Input:} P1=$(2,3)$, P2=$(2,1)$, P3=$(8,2)$, P4=$(9,1)$.

\textbf{Hand decision trace.} The CPU is idle until time 2. P1 and P2 arrive
together, so earlier input position makes P1 first. The CPU becomes idle again at
time 6. P3 begins at time 8 and P4 waits until P3 completes.

\timeline{0--2 IDLE \;|\; 2--5 P1 \;|\; 5--6 P2 \;|\; 6--8 IDLE \;|\; 8--10 P3 \;|\; 10--11 P4}

\begin{table}[H]
  \centering\small
  \caption{FCFS Workload B: hand calculation and simulator output (identical).}
  \begin{tabular}{lrrrrrrr}
    \toprule
    Process                              & A             & B             & FS            & C  & T & W & R \\
    \midrule
    P1                                   & 2             & 3             & 2             & 5  & 3 & 0 & 0 \\
    P2                                   & 2             & 1             & 5             & 6  & 4 & 3 & 3 \\
    P3                                   & 8             & 2             & 8             & 10 & 2 & 0 & 0 \\
    P4                                   & 9             & 1             & 10            & 11 & 2 & 1 & 1 \\
    \midrule
    \multicolumn{5}{r}{\textbf{Average}} & \textbf{2.75} & \textbf{1.00} & \textbf{1.00}                  \\
    \bottomrule
  \end{tabular}
\end{table}

\textbf{Simulator comparison.} The simulator produced the same running and idle
intervals and every per-process value in the table. \exactmatch

\simshot{b-fcfs.png}{Simulator output for FCFS on Workload B.}

\subsection{SJF on Workload B}

\textbf{Input:} P1=$(2,3)$, P2=$(2,1)$, P3=$(8,2)$, P4=$(9,1)$.

\textbf{Hand decision trace.} At time 2, P2 has the shortest burst, then P1 runs.
At time 8 only P3 is ready, and non-preemptive SJF does not interrupt it when P4
arrives.

\timeline{0--2 IDLE \;|\; 2--3 P2 \;|\; 3--6 P1 \;|\; 6--8 IDLE \;|\; 8--10 P3 \;|\; 10--11 P4}

\begin{table}[H]
  \centering\small
  \caption{SJF Workload B: hand calculation and simulator output (identical).}
  \begin{tabular}{lrrrrrrr}
    \toprule
    Process                              & A             & B             & FS            & C  & T & W & R \\
    \midrule
    P1                                   & 2             & 3             & 3             & 6  & 4 & 1 & 1 \\
    P2                                   & 2             & 1             & 2             & 3  & 1 & 0 & 0 \\
    P3                                   & 8             & 2             & 8             & 10 & 2 & 0 & 0 \\
    P4                                   & 9             & 1             & 10            & 11 & 2 & 1 & 1 \\
    \midrule
    \multicolumn{5}{r}{\textbf{Average}} & \textbf{2.25} & \textbf{0.50} & \textbf{0.50}                  \\
    \bottomrule
  \end{tabular}
\end{table}

\textbf{Simulator comparison.} The simulator produced the same timeline and every
per-process value in the table. \exactmatch

\simshot{b-sjf.png}{Simulator output for SJF on Workload B.}

\subsection{SRTF on Workload B}

\textbf{Input:} P1=$(2,3)$, P2=$(2,1)$, P3=$(8,2)$, P4=$(9,1)$.

\textbf{Hand decision trace.} P2 runs before P1 at time 2. At time 9, P3 has one
unit remaining and the newly arrived P4 also requires one unit. Equal remaining
time does not cause preemption, so P3 continues. The schedule is identical to SJF.

\timeline{0--2 IDLE \;|\; 2--3 P2 \;|\; 3--6 P1 \;|\; 6--8 IDLE \;|\; 8--10 P3 \;|\; 10--11 P4}

\begin{table}[H]
  \centering\small
  \caption{SRTF Workload B: hand calculation and simulator output (identical).}
  \begin{tabular}{lrrrrrrr}
    \toprule
    Process                              & A             & B             & FS            & C  & T & W & R \\
    \midrule
    P1                                   & 2             & 3             & 3             & 6  & 4 & 1 & 1 \\
    P2                                   & 2             & 1             & 2             & 3  & 1 & 0 & 0 \\
    P3                                   & 8             & 2             & 8             & 10 & 2 & 0 & 0 \\
    P4                                   & 9             & 1             & 10            & 11 & 2 & 1 & 1 \\
    \midrule
    \multicolumn{5}{r}{\textbf{Average}} & \textbf{2.25} & \textbf{0.50} & \textbf{0.50}                  \\
    \bottomrule
  \end{tabular}
\end{table}

\textbf{Simulator comparison.} The simulator produced the same timeline and every
per-process value in the table. \exactmatch

\simshot{b-srtf.png}{Simulator output for SRTF on Workload B.}

\subsection{Round Robin (q=2) on Workload B}

\textbf{Input:} P1=$(2,3)$, P2=$(2,1)$, P3=$(8,2)$, P4=$(9,1)$; quantum $q=2$.

\textbf{Hand decision trace.} P1 and P2 arrive together, so input order places P1
first. After its two-unit slice, P2 runs before re-enqueued P1. P3 completes its
two-unit slice even though P4 arrives at time 9.

\timeline{0--2 IDLE \;|\; 2--4 P1 \;|\; 4--5 P2 \;|\; 5--6 P1 \;|\; 6--8 IDLE \;|\; 8--10 P3 \;|\; 10--11 P4}

\begin{table}[H]
  \centering\small
  \caption{Round Robin Workload B: hand calculation and simulator output (identical).}
  \begin{tabular}{lrrrrrrr}
    \toprule
    Process                              & A             & B             & FS            & C  & T & W & R \\
    \midrule
    P1                                   & 2             & 3             & 2             & 6  & 4 & 1 & 0 \\
    P2                                   & 2             & 1             & 4             & 5  & 3 & 2 & 2 \\
    P3                                   & 8             & 2             & 8             & 10 & 2 & 0 & 0 \\
    P4                                   & 9             & 1             & 10            & 11 & 2 & 1 & 1 \\
    \midrule
    \multicolumn{5}{r}{\textbf{Average}} & \textbf{2.75} & \textbf{1.00} & \textbf{0.75}                  \\
    \bottomrule
  \end{tabular}
\end{table}

\textbf{Simulator comparison.} The simulator produced the same timeline and every
per-process value in the table. \exactmatch

\simshot{b-rr.png}{Simulator output for Round Robin ($q=2$) on Workload B.}

\section{Additional test cases}

\subsection{Idle CPU test}

\textbf{Input:} A=$(2,2)$ and B=$(6,1)$; algorithm FCFS.

\textbf{Behavior checked:} The initial and middle idle intervals must be visible,
and the CPU must restart as soon as the next process arrives.

\timeline{Expected: 0--2 IDLE \;|\; 2--4 A \;|\; 4--6 IDLE \;|\; 6--7 B}

Expected A: $FS=2$, $C=4$, $T=2$, $W=0$, $R=0$. Expected B: $FS=6$,
$C=7$, $T=1$, $W=0$, $R=0$. The expected averages are $T=1.50$, $W=0$,
and $R=0$. The simulator matched these results.

\subsection{Simultaneous arrivals and scheduling tie}

\textbf{Input:} P1=$(0,2)$, P2=$(0,2)$, P3=$(1,1)$; algorithm SJF.

\textbf{Behavior checked:} P1 and P2 tie on arrival and burst, so earlier input
position selects P1. When it finishes, the shorter P3 runs before P2.

\timeline{Expected: 0--2 P1 \;|\; 2--3 P3 \;|\; 3--5 P2}

Expected $(T,W,R)$: P1=$(2,0,0)$, P2=$(5,3,3)$, P3=$(2,1,1)$. The expected
averages are $T=3.00$, $W=1.33$, and $R=1.33$. The simulator matched these results.

\subsection{SRTF preemption and equal-remaining rule}

\textbf{Input:} P1=$(0,7)$, P2=$(2,2)$, P3=$(3,1)$; algorithm SRTF.

\textbf{Behavior checked:} P2 strictly preempts P1 at time 2. At time 3, P2 and
newly arrived P3 both have one unit remaining, so P2 must continue rather than be
preempted.

\timeline{Expected: 0--2 P1 \;|\; 2--4 P2 \;|\; 4--5 P3 \;|\; 5--10 P1}

Expected $(T,W,R)$: P1=$(10,3,0)$, P2=$(2,0,0)$, P3=$(2,1,1)$. The expected
averages are $T=4.67$, $W=1.33$, and $R=0.33$. The simulator matched these results.

\subsection{Verified schedule properties}

\begin{table}[H]
  \centering\small
  \begin{tabularx}{\linewidth}{>{\bfseries}lX>{\color{success}\bfseries}l}
    \toprule
    Property                    & Verification method                                                           & Result \\
    \midrule
    Exact CPU burst             & Sum every process interval and compare it with the requested burst.           & Pass   \\
    No execution before arrival & Require every process interval start to be at or after its arrival.           & Pass   \\
    No CPU overlap              & Require each interval to start exactly when the previous interval ends.       & Pass   \\
    Nonnegative metrics         & Check every waiting and response value.                                       & Pass   \\
    Work conserving             & For every idle interval, confirm that no arrived process has unfinished work. & Pass   \\
    \bottomrule
  \end{tabularx}
\end{table}

The automated test suite applies these invariants to all four algorithms on both
required workloads. Exact-output tests also cover ties, Round Robin boundary
ordering, preserved RR slices, and rejected invalid input.

\section{Analysis}

\subsection{Which algorithm is best for Workload A?}

\begin{table}[H]
  \centering
  \begin{tabular}{lrrr}
    \toprule
    Algorithm  & Average T & Average W     & Average R     \\
    \midrule
    FCFS       & 10.50     & 6.75          & 6.75          \\
    SJF        & 9.00      & 5.25          & 5.25          \\
    SRTF       & 6.25      & \textbf{2.50} & \textbf{0.00} \\
    RR ($q=2$) & 8.50      & 4.75          & 1.75          \\
    \bottomrule
  \end{tabular}
\end{table}

SRTF gives both the lowest average waiting time (2.50) and the lowest average
response time (0.00). Each later job in A is shorter than the older job's remaining
time when it first runs, so it preempts and receives an immediate first response.
This favors the short arrivals while P1 absorbs most of the waiting. SJF knows the
burst lengths but cannot interrupt P1's initial eight-unit burst.

\subsection{Round Robin with quantum 1, 2, and 4}

\begin{table}[H]
  \centering
  \begin{tabular}{rrrrr}
    \toprule
    Quantum & Slices & Average T & Average W & Average R \\
    \midrule
    1       & 15     & 8.50      & 4.75      & 0.75      \\
    2       & 8      & 8.50      & 4.75      & 1.75      \\
    4       & 5      & 9.25      & 5.50      & 3.75      \\
    \bottomrule
  \end{tabular}
\end{table}

For Workload A, $q=1$ improves initial responsiveness but does not improve average
turnaround or waiting relative to $q=2$. A smaller quantum moves first starts
earlier, but it also interrupts jobs more often and may postpone completion through
additional queue rotations. This simulator assumes zero context-switch cost. On a
real CPU, an extremely small quantum would additionally lose useful time to context
switching. Therefore, decreasing the quantum is not guaranteed to improve every
metric.

\subsection{When SJF and SRTF produce the same schedule}

SJF and SRTF produce the same schedule when, after each dispatch, no later arrival
has a burst strictly shorter than the running process's remaining time. This includes
a batch where all jobs are already available, or later arrivals that cannot justify
preemption. Workload B is an example. At time 9, P4 arrives while P3 has one unit
remaining. P4 also needs one unit, so SRTF's equality rule keeps P3 running and the
SJF and SRTF schedules remain identical.

\subsection{Why real operating systems estimate CPU bursts}

The simulator receives the complete synthetic workload in advance, including the
exact duration of every future CPU burst. A real operating system cannot know how
long a program's next CPU burst will last because execution depends on inputs,
branches, cache behavior, I/O, and future events. A practical scheduler must use
estimates from observed history or indirect priority and responsiveness signals.
This simulator can make ideal SJF/SRTF decisions only because the input gives it
perfect future knowledge.

\subsection{Starvation under SJF and SRTF}

Starvation can occur under both policies. Under SJF, a ready long job may repeatedly
lose dispatch opportunities if shorter jobs keep arriving before the CPU becomes
free. Under SRTF, a long job may be delayed or preempted whenever strictly shorter
work arrives. The finite workloads in this assignment eventually stop producing
new jobs, so every process finishes. This removes the unbounded-arrival condition
needed for indefinite starvation. Aging is one possible mitigation in a real
scheduler.

\subsection{Implementation mistake exposed by testing}

The first browser render showed the Round Robin quantum field even while FCFS was
selected. The HTML correctly used the \texttt{hidden} attribute, but the component
CSS rule set \texttt{display:grid} and overrode the browser's default hidden style.
A headless-browser screenshot exposed the problem. I added an explicit
\texttt{[hidden] \{ display:none !important; \}} rule and repeated the visual check.
The scheduling-engine tests alone would not have detected this UI-only error.

\section{Implementation notes}

The scheduling engine is a pure JavaScript module separated from the page renderer.
FCFS and SJF dispatch whole jobs. SRTF advances one integer time unit so arrivals
and the strict-preemption rule are evaluated exactly. Round Robin maintains an
explicit FIFO queue, processes arrivals through each boundary, and only then
re-enqueues an unfinished runner. Non-RR adjacent intervals are merged for chart
readability; RR slices are intentionally kept separate.

Invalid input produces specific messages for missing or duplicate IDs, invalid
arrival/burst values, unsupported algorithms, and an invalid Round Robin quantum.
The site requires no backend and stores no personal data. It models the assignment
policies and does not reproduce Ubuntu's actual scheduler.

\section{AI assistance statement}

I used OpenAI Codex to help implement the HTML/CSS/JavaScript interface and pure
scheduling module, draft automated tests, check arithmetic, debug the UI, and
organize this report. I checked the generated work by independently tracing
Workloads A and B, comparing every timeline and metric, running exact-output and
invariant tests, and rendering the interface in a headless browser.

I rejected one generated-code optimization for Round Robin: the general timeline
helper normally merges adjacent intervals belonging to the same process. Doing so
would hide a quantum boundary when no other process is ready, contrary to the
assignment. Round Robin therefore records consecutive slices separately.

\end{document}
